%=====================================================================
\chapter{SQL: Creating the Schema}\label{ch:ddl}
%=====================================================================
\locator{3}{Part III --- From Diagram to Schema}

\begin{outcomes}
By the end of this chapter you will be able to:
\begin{tight}
  \item \textbf{Write} a complete \sqlkw{CREATE TABLE} statement with column
        types, key, nullability, defaults and \sqlkw{CHECK} constraints.
  \item \textbf{Explain} what type affinity is, and \textbf{predict} what
        SQLite stores for a given value in a given declared type.
  \item \textbf{Demonstrate} that an \texttt{INTEGER} column will hold text,
        and that \sqlkw{STRICT} stops it.
  \item \textbf{Trace} a wrong query result back to a value of the wrong type.
  \item \textbf{Choose} between a column constraint and a table constraint, and
        \textbf{name} a case where only the second will do.
\end{tight}
\end{outcomes}

\begin{prereq}
\chref{mapping}: the tables and keys you are now declaring.
\chref{constraints} for \sqlkw{PRAGMA} \texttt{foreign\_keys} and \sqlkw{NOT
NULL}, both of which appear again here for a second reason. You need
\texttt{code/m10\_ddl.py}.
\end{prereq}

\begin{keyterms}
DDL $\bullet$ \sqlkw{CREATE TABLE} $\bullet$ \sqlkw{DROP TABLE} $\bullet$
storage class $\bullet$ type affinity $\bullet$ declared type $\bullet$
\sqlkw{STRICT} $\bullet$ \texttt{typeof()} $\bullet$ column constraint
$\bullet$ table constraint $\bullet$ \sqlkw{CHECK} $\bullet$ \sqlkw{DEFAULT}
$\bullet$ \sqlkw{UNIQUE} $\bullet$ \sqlkw{AUTOINCREMENT} $\bullet$ rowid
$\bullet$ generated column
\end{keyterms}

\section{The Statement, in Full}

\sqlkw{CREATE TABLE} is the whole of data definition that this course needs.
Everything in it is a promise, and the order of the clauses is fixed.

\begin{definitionbox}[The anatomy of a \sqlkw{CREATE TABLE}]
\begin{lstlisting}[language=sql]
CREATE TABLE student (
  roll_no   TEXT    PRIMARY KEY NOT NULL,     -- column constraints
  name      TEXT    NOT NULL,
  programme TEXT    NOT NULL DEFAULT 'BSIT',
  semester  INTEGER NOT NULL CHECK (semester BETWEEN 1 AND 8),
  cnic      TEXT    UNIQUE,
  CHECK (length(roll_no) = 11)                -- a table constraint
) STRICT;
\end{lstlisting}

\smallskip
A \term{column constraint} follows its column and can mention only that
column. A \term{table constraint} is listed after all the columns and may
mention several. \sqlkw{PRIMARY KEY} on more than one column, and any
\sqlkw{CHECK} comparing two columns, must be table constraints --- there is no
other way to write them.
\end{definitionbox}

\begin{examplebox}[The five clauses, and what each one is for]
\begin{tabular}{@{}L{3.0cm}L{4.5cm}L{6.5cm}@{}}
\textbf{Clause} & \textbf{Promises} & \textbf{Comes from} \\
\sqlkw{PRIMARY KEY} & these values identify a row & the key chosen in
  \chref{ermodel} \\
\sqlkw{NOT NULL} & a value is always present & total participation
  (\chref{cardinality}), or entity integrity \\
\sqlkw{UNIQUE} & no two rows agree here & an \emph{alternate} key \\
\sqlkw{CHECK} & every value satisfies this & the attribute's domain
  (\chref{constraints}) \\
\sqlkw{DEFAULT} & what to use if none given & a university policy, not a
  guess \\
\end{tabular}

\smallskip
\textbf{\sqlkw{DEFAULT} is the one to be careful with.} It turns a missing
value into a real one, so a row that \emph{should} have been rejected is
accepted with a plausible value instead. Use it where the university genuinely
has a default --- \texttt{semester} is 1 for a new admission --- and not to
avoid thinking about nulls.
\end{examplebox}

\section{Type Affinity, and What It Costs}

Every other relational system treats a declared type as a rule. SQLite treats
it as a \emph{preference}, and the difference is the subject of this chapter.

\begin{definitionbox}[Storage class, declared type, affinity]
A \term{storage class} is what a value actually \emph{is} in the file: null,
integer, real, text or blob. Every value has one, and \texttt{typeof()} reports
it.

\smallskip
A \term{declared type} is what you wrote in \sqlkw{CREATE TABLE}.

\smallskip
\term{Type affinity} is the rule connecting them: on insert, SQLite
\emph{tries} to convert the value to the column's declared type. If it can, it
does. \textbf{If it cannot, it stores the value unchanged.} No error.

\smallskip
\sqlkw{STRICT}, added to SQLite in 2021, changes the last sentence: the value
is refused. A \sqlkw{STRICT} table may only use the types \texttt{INT},
\texttt{INTEGER}, \texttt{REAL}, \texttt{TEXT}, \texttt{BLOB} and
\texttt{ANY}.
\end{definitionbox}

\dsfig{%
\begin{tikzpicture}[
  val/.style={rounded corners=2pt, draw=secondary, line width=0.9pt,
              fill=secondary!8, text=secondary!60!black,
              font=\scriptsize\ttfamily, minimum width=1.7cm,
              minimum height=0.6cm},
  dec/.style={draw=primary, line width=1pt, fill=primary!7,
              text=primary!60!black, font=\scriptsize\bfseries,
              align=center, text width=3.1cm, minimum height=1.0cm,
              rounded corners=3pt},
  good/.style={rounded corners=2pt, draw=greenhl, line width=0.9pt,
               fill=greenhl!8, text=greenhl!55!black,
               font=\scriptsize\bfseries, align=center, text width=3.0cm,
               minimum height=0.75cm},
  bad/.style={good, draw=accent, fill=accent!8, text=accent!75!black},
  strict/.style={good, draw=purplehl, fill=purplehl!8,
                 text=purplehl!65!black},
  arr/.style={-{Stealth[length=2.2mm]}, primary!55, line width=0.9pt},
  tag/.style={font=\scriptsize\bfseries, text=primary!70!black},
  nt/.style={font=\scriptsize\itshape, text=neutral, align=left}
]
  \node[val] (v) at (-6.0,0) {'hello'};
  \node[nt, anchor=south, text width=2.6cm, align=center]
       at (-6.0,0.45) {a value, on its way into\\\texttt{c INTEGER}};

  \node[dec] (d) at (-2.4,0) {can it be converted\\to the column's\\declared type?};
  \draw[arr] (v) -- (d);

  \node[good]   (yes) at ( 2.0, 1.5) {converted\\{\scriptsize\mdseries
       \texttt{'42'} $\rightarrow$ integer 42}};
  \node[bad]    (no)  at ( 2.0, 0.0) {\textbf{stored unchanged}\\{\scriptsize
       \mdseries \texttt{'hello'} stays text}};
  \node[strict] (st)  at ( 2.0,-1.6) {refused\\{\scriptsize\mdseries
       \texttt{cannot store TEXT}}};

  \draw[arr] (d) -- node[tag, above, pos=0.6] {yes} (yes);
  \draw[arr] (d) -- node[tag, above, pos=0.62] {no} (no);
  % pos=0.33, not 0.6: further along the line the label landed on the box it
  % points at, and the arrowhead disappeared behind its own caption.
  \draw[arr, purplehl!60] (d) -- node[tag, below, pos=0.33,
       text=purplehl!70!black] {no, and \sqlkw{STRICT}} (st);

  \node[nt, anchor=west, text width=4.2cm, text=greenhl!55!black]
       at (3.9, 1.5) {silently succeeds};
  \node[nt, anchor=west, text width=4.4cm, text=accent!75!black]
       at (3.9, 0.0) {\textbf{silently fails.} No error, and the column
        looks the same as the row above it};
  \node[nt, anchor=west, text width=4.2cm, text=purplehl!65!black]
       at (3.9,-1.6) {the only outcome you can reason about};
\end{tikzpicture}}%
{Type affinity in one decision. The middle branch is the whole problem: it is
not an error, it is not reported, and a column containing both of the top two
outcomes looks exactly like a column containing only the first.}{affinity}

\begin{measured}[A column declared \texttt{INTEGER} that holds the word ``absent'']
\textsf{\footnotesize \texttt{code/m10\_ddl.py}, SQLite 3.51.2, Intel i5-6600,
16\,GB, Windows 11.}

\medskip
\textbf{1. One value, five declared types --- what is actually stored}

\begin{tabular}{@{}L{2.8cm}L{2.0cm}L{1.8cm}L{1.8cm}L{2.2cm}L{2.2cm}@{}}
\textbf{inserted} & \texttt{INTEGER} & \texttt{TEXT} & \texttt{REAL} &
\texttt{NUMERIC} & \textbf{(no type)} \\
\texttt{42} & integer & text & real & integer & integer \\
\texttt{'42'} & \textbf{integer} & text & real & integer & text \\
\texttt{'hello'} & \textbf{text} & text & \textbf{text} & \textbf{text} & text \\
\texttt{3.7} & \textbf{real} & text & real & real & real \\
\sqlkw{NULL} & null & null & null & null & null \\
\end{tabular}

\smallskip
Read the \texttt{'hello'} row. Five columns, five different declared types,
and every one of them stored the string.

\medskip
\textbf{2. The same column, one keyword different}

\begin{tabular}{@{}L{3.4cm}L{4.4cm}L{6.0cm}@{}}
\textbf{inserted into \texttt{c INTEGER}} & \textbf{plain} & \textbf{\sqlkw{STRICT}} \\
\texttt{42}      & stored as integer & stored as integer \\
\texttt{'42'}    & stored as integer & stored as integer \\
\texttt{'hello'} & \textbf{stored as text} & refused: \texttt{cannot store TEXT value in INT} \\
\texttt{3.7}     & \textbf{stored as real} & refused: \texttt{cannot store REAL value in INT} \\
\sqlkw{NULL}     & stored as null & stored as null \\
\end{tabular}

\medskip
\textbf{3. What it does to the answers}

Marks entered as \texttt{A=95, B=72, C='absent', D=88, E='100'}:

\begin{tabular}{@{}L{5.6cm}R{3.2cm}L{5.4cm}@{}}
\sqlkw{ORDER BY} \texttt{score} & \texttt{absent} last &
  text sorts \emph{above} every number \\
\sqlkw{MAX}(\texttt{score}) & \texttt{'absent'} &
  so C is top of the class \\
\sqlkw{AVG}(\texttt{score}) as written & \textbf{71.00} &
  \texttt{'absent'} read as 0, still divided by 5 \\
\sqlkw{AVG} over the numeric rows & 88.75 & the honest figure \\
\textbf{the reported average is} & \textbf{20\% out} & and nothing says so \\
\end{tabular}

\smallskip
Note also that \texttt{'100'} \emph{was} converted, to the integer 100.
Affinity succeeded silently there and failed silently on \texttt{'absent'},
and the column looks identical either way.

\medskip
\textbf{4. What \sqlkw{STRICT} costs} (50{,}000 rows, 5 runs)

\begin{tabular}{@{}L{5.0cm}R{2.6cm}R{2.4cm}R{2.4cm}@{}}
& \textbf{median} & \textbf{min} & \textbf{max} \\
plain          & 67.9\,ms & 66.2 & 78.5 \\
\sqlkw{STRICT} & 67.9\,ms & 66.8 & 72.8 \\
\textbf{difference} & \textbf{0.0\,ms} & \multicolumn{2}{l}{\textbf{(+0\%)}} \\
\end{tabular}

\smallskip
\textbf{The two ranges overlap}, so the difference is noise rather than a
cost --- which is the opposite of \chref{constraints}'s foreign-key figure,
where they did not overlap and the 24\% was real. \sqlkw{STRICT} is free: it
is checking a type it already knows.

\smallskip
\textbf{So the whole argument is one-sided.} The protection costs nothing
measurable and the absence of it silently moved a class average by twenty per
cent. Every table in this course is \sqlkw{STRICT}.
\end{measured}

\begin{alertbox}[Why SQLite is like this, and why that does not help you]
This is not a bug and it will not be fixed. SQLite was designed to be
embedded, tolerant and backward-compatible for decades; rejecting a value that
an older version accepted would break databases in the field. Affinity is the
documented, deliberate behaviour, and \sqlkw{STRICT} was added as an opt-in
twenty-one years later.

\smallskip
Knowing the reason does not change what you must do:

\begin{tight}
  \item \textbf{Write \sqlkw{STRICT} on every table.} There is no measured
        reason not to.
  \item \textbf{Use \texttt{typeof()} when a result looks wrong.} It is the
        first diagnostic, and the one nobody reaches for.
  \item \textbf{Be most suspicious of data you imported.} Values typed by hand
        into \sqlkw{INSERT} statements are usually the right type; values read
        from a CSV are all text until something converts them, and
        \chref{evolution} is where that bites.
  \item \textbf{Do not assume other systems behave this way.} PostgreSQL and
        MySQL reject \texttt{'hello'} in an integer column outright. This
        chapter's behaviour is SQLite's, and \chref{project} is where the
        comparison is drawn properly.
\end{tight}
\end{alertbox}

\section{Three Smaller Things That Matter}

\begin{conceptbox}[Rowid, \sqlkw{AUTOINCREMENT}, and generated columns]
\textbf{Every ordinary SQLite table has a hidden \texttt{rowid}} --- a 64-bit
integer key the system maintains. Declaring a column \texttt{INTEGER PRIMARY
KEY} makes that column \emph{be} the rowid rather than sit beside it, which is
why \chref{constraints} measured it substituting an integer for a null.

\smallskip
\textbf{\sqlkw{AUTOINCREMENT} is almost never what you want.} Without it,
SQLite reuses the largest rowid when rows are deleted; with it, it never does,
at the cost of an extra table and a write per insert. Use it only when a
reused identifier would be a real problem --- an invoice number, not a
\texttt{section\_id}.

\smallskip
\textbf{A generated column} is computed from the others and never stored
wrongly, which makes it the right home for \chref{ermodel}'s derived
attributes:

\begin{lstlisting}[language=sql]
  date_of_birth TEXT NOT NULL,
  age INTEGER GENERATED ALWAYS AS
      (CAST((julianday('now') - julianday(date_of_birth)) / 365.25
            AS INTEGER)) VIRTUAL
\end{lstlisting}

\sqlkw{VIRTUAL} computes it on every read and stores nothing;
\sqlkw{STORED} writes it to disk. For \texttt{age}, \sqlkw{VIRTUAL} is the
only correct choice: a stored one would be wrong by morning, which was exactly
\chref{ermodel}'s objection to storing it at all.
\end{conceptbox}

\begin{pitfall}
\begin{tight}
  \item \textbf{Trusting a declared type.} Without \sqlkw{STRICT} it is a
        preference. Measured above.
  \item \textbf{Using \texttt{VARCHAR(50)}.} SQLite accepts the words and
        ignores the 50 entirely. Write \texttt{TEXT}, and enforce a length
        with \sqlkw{CHECK} if you mean one.
  \item \textbf{Declaring a composite primary key as two column
        constraints.} \texttt{a INTEGER PRIMARY KEY, b TEXT PRIMARY KEY} is an
        error; it must be one table constraint.
  \item \textbf{Using \sqlkw{DEFAULT} to dodge a null.} It replaces a
        rejection with a plausible wrong value.
  \item \textbf{Reaching for \sqlkw{AUTOINCREMENT} by habit.} It costs a write
        per insert and solves a problem you probably do not have.
  \item \textbf{Storing a derived value.} Use a \sqlkw{VIRTUAL} generated
        column.
  \item \textbf{Forgetting that \sqlkw{CHECK} is not re-run on existing rows.}
        \chref{evolution} is about exactly that.
\end{tight}
\end{pitfall}

\begin{lab}[Store the wrong type on purpose, then make it impossible]
\begin{lstlisting}[language=bash]
# --- 1. A scratch database; keep university.db clean ----------------
sqlite3 types.db
\end{lstlisting}
\begin{lstlisting}[language=sql]
-- --- 2. A table that looks completely ordinary ------------------
.headers on
.mode box
CREATE TABLE marks (roll TEXT PRIMARY KEY, score INTEGER NOT NULL);

-- --- 3. Put four marks in, one of them the word 'absent' --------
INSERT INTO marks VALUES ('A',95),('B',72),('C','absent'),('D',88);
SELECT roll, score, typeof(score) FROM marks;
-- One row is text. Nothing warned you. The column says INTEGER.

-- --- 4. Ask three ordinary questions ----------------------------
SELECT roll, score FROM marks ORDER BY score DESC;
SELECT MAX(score) FROM marks;
SELECT AVG(score) FROM marks;
-- Write all three answers down. Then work out, by hand, what each
-- one SHOULD be. Two of the three are wrong.

-- --- 5. Find the bad row, which is the real skill ---------------
SELECT roll, score FROM marks
 WHERE typeof(score) NOT IN ('integer','real');
-- Learn this query. It is how you audit a column you inherited.

-- --- 6. The honest average --------------------------------------
SELECT AVG(score) FROM marks
 WHERE typeof(score) IN ('integer','real');
-- Compare with step 4. How far out was it, as a percentage?

-- --- 7. Make it impossible --------------------------------------
CREATE TABLE marks2 (
  roll  TEXT    PRIMARY KEY NOT NULL,
  score INTEGER NOT NULL CHECK (score BETWEEN 0 AND 100)
) STRICT;
INSERT INTO marks2 VALUES ('A',95);
INSERT INTO marks2 VALUES ('C','absent');
INSERT INTO marks2 VALUES ('E',150);
-- Two refusals, two different messages. Record both, and say
-- which constraint produced each.

-- --- 8. Prove VARCHAR(50) is not a length limit -----------------
CREATE TABLE v (s VARCHAR(5));
INSERT INTO v VALUES ('this string is far longer than five');
SELECT length(s) FROM v;
-- Say in one sentence what VARCHAR(5) accomplished.
\end{lstlisting}
\begin{lstlisting}[language=bash]
# --- 9. The measured version ----------------------------------------
python code/m10_ddl.py

# --- 10. Apply it to your own schema ---------------------------------
# Every table in code/schema/v1_create.sql must end in STRICT.
# Check, and then prove it:
sqlite3 university.db "SELECT name, sql FROM sqlite_master
  WHERE type='table' AND sql NOT LIKE '%STRICT%';"
# That query should return nothing.
\end{lstlisting}

\smallskip
\textbf{What to notice.} Step 4 is the point of the chapter: three ordinary
questions, no errors, two wrong answers. Step 5 is the skill --- a
\texttt{typeof()} audit is the only way to find this in a database you did not
build, and it is worth running on any data you import. Step 10 is the
protection, and it costs nothing measurable.
\end{lab}

\begin{checkpoint}
\begin{tightnum}
  \item A column is declared \texttt{INTEGER} and \sqlkw{AVG} over it returns
        71 when the true average of the real marks is 88.75. Explain exactly
        how the database arrived at 71, and give the two queries you would run
        to prove your explanation.
  \item State the difference between a column constraint and a table
        constraint, and give two constraints from your own
        \texttt{v1\_create.sql} that \emph{must} be table constraints, saying
        why each one cannot be written as a column constraint.
  \item \sqlkw{STRICT} costs nothing measurable and prevents a class of silent
        error. Give the one argument somebody could still make against putting
        it on every table, and say whether you accept it.
\end{tightnum}
\end{checkpoint}

\begin{chaptersummary}
\begin{tight}
  \item \sqlkw{CREATE TABLE} carries five kinds of promise: \sqlkw{PRIMARY
        KEY}, \sqlkw{NOT NULL}, \sqlkw{UNIQUE}, \sqlkw{CHECK} and
        \sqlkw{DEFAULT}, each traceable to a decision in Part~II.
  \item A \textbf{column constraint} mentions one column; a \textbf{table
        constraint} may mention several. Composite keys and cross-column
        \sqlkw{CHECK}s must be table constraints.
  \item A \textbf{storage class} is what a value is; a \textbf{declared type}
        is what you asked for; \textbf{affinity} converts if it can and
        \textbf{stores the value unchanged if it cannot}.
  \item \textbf{Measured:} \texttt{'hello'} went into columns declared
        \texttt{INTEGER}, \texttt{REAL} and \texttt{NUMERIC}, unchanged and
        without an error.
  \item \textbf{Measured:} on a mixed column, \sqlkw{ORDER BY} sorts text above
        every number, \sqlkw{MAX} returns \texttt{'absent'}, and \sqlkw{AVG}
        reads it as 0 while still dividing by the full count --- a class
        average reported \textbf{20\% low}, silently.
  \item \texttt{'100'} \emph{was} converted to 100. Affinity succeeds and fails
        silently, and the column looks the same either way.
  \item \textbf{Measured: \sqlkw{STRICT} costs nothing} --- the two timing
        ranges overlap, unlike \chref{constraints}'s foreign-key figure. The
        argument for it is one-sided.
  \item \texttt{typeof()} is the first diagnostic when a result looks wrong,
        and the one nobody reaches for. Be most suspicious of imported data.
  \item \texttt{VARCHAR(50)} has no length limit in SQLite. Use \texttt{TEXT}
        plus a \sqlkw{CHECK}.
  \item A \textbf{\sqlkw{VIRTUAL} generated column} is the right home for a
        derived attribute: computed on read, never stale.
\end{tight}
\end{chaptersummary}

\begin{reviewq}
  \item \textbf{(MCQ)} Inserting \texttt{'hello'} into a column declared
        \texttt{INTEGER} in a non-\sqlkw{STRICT} table: (a)~raises an error;
        (b)~stores 0; (c)~stores the text unchanged; (d)~stores null.
  \item \textbf{(MCQ)} \texttt{VARCHAR(5)} in SQLite limits the column to:
        (a)~5 characters; (b)~5 bytes; (c)~nothing; (d)~5 characters only in
        \sqlkw{STRICT} tables.
  \item \textbf{(Short)} Why does \texttt{'42'} become an integer while
        \texttt{'hello'} stays text? State the rule in one sentence.
  \item \textbf{(Short)} Give a constraint that can only be written as a table
        constraint, and say why.
  \item \textbf{(Short)} \sqlkw{STRICT} showed no measurable cost while foreign
        key enforcement showed 24\%. Explain why the two results differ.
  \item \textbf{(Applied)} You are given a table of 50{,}000 imported rows with
        columns \texttt{roll\_no}, \texttt{score} and \texttt{date\_sat}.
        Write the queries that would tell you whether any value is of the wrong
        type, decide which results you would report as faults, and describe
        how you would repair the table without losing a row.
  \item \textbf{(Applied)} Rewrite one table of your \texttt{v1\_create.sql} to
        add a \sqlkw{DEFAULT}, a cross-column \sqlkw{CHECK} and a
        \sqlkw{VIRTUAL} generated column. For each, say what fact about the
        university it records and what would go wrong if it were omitted ---
        and for the \sqlkw{DEFAULT}, say what it would wrongly accept.
\end{reviewq}
